\documentclass{article}
\usepackage[UTF8]{ctex}
\usepackage{amsmath,amssymb,amsfonts}
\usepackage{enumitem}
\usepackage{geometry}
\geometry{a4paper, margin=1in}
% === 错题本排版优化（不参与题目标识） ===
\setlength{\parindent}{0pt}
\setlength{\parskip}{0.35em plus 0.10em minus 0.08em}
\setlength{\emergencystretch}{3em}
\setlist{itemsep=0.25em, topsep=0.40em, parsep=0pt}
\allowdisplaybreaks[2]
% === 错题本排版优化结束 ===
\begin{document}
\section*{北京师范大学2026年数学分析真题}
\section*{3．北京师范大学2026年数学分析试卷}

\subsection{题目1}
一．（15 分）证明：

\[
1-\frac{1}{2}+\frac{1}{3}-\frac{1}{4}+\cdots+\frac{1}{2025}-\frac{1}{2026}<\frac{\sqrt{2}}{2} .
\]

\subsection{题目2}
二．（15 分）设 \(f: \mathbb{R} \times \mathbb{R} \rightarrow \mathbb{R}\) 满足对任意的 \(a \in \mathbb{R}\) ，有 \(f(0, a) \leq f(0,0) \leq f(a, 0)\) ．证明：

\[
\sup _{b \in \mathbb{R}} \inf _{a \in \mathbb{R}} f(a, b)=\inf _{a \in \mathbb{R}} \sup _{b \in \mathbb{R}} f(a, b) .
\]

\subsection{题目3}
三．（15 分）求极限 \(\lim _{x \rightarrow+\infty} x^{2}\left[\left(\frac{3 x^{3}-7}{3 x^{3}-x}\right)^{\frac{1}{5}}-x \sin \frac{1}{x}\right]\) ．

\subsection{题目4}
四．（15 分）已知 \(f(x)=\sum_{n=1}^{\infty} \frac{\cos (n x)}{\sqrt{n^{3}+n}}\) ．证明：

\begin{itemize}
  \item[（1）] （7 分）\(f(x)\) 在 \((-\infty,+\infty)\) 上连续．
  \item[（2）] （8 分）\(\frac{\sqrt{2}}{2}-\frac{1}{15}<\int_{0}^{\frac{\pi}{2}} f(x) \mathrm{d} x<\frac{\sqrt{2}}{2}\) ．
\end{itemize}

\subsection{题目5}
五．（15 分）设 \(f(x)\) 在 \([1,+\infty)\) 上非负单调递减，证明： \(\lim _{n \rightarrow \infty}\left[\sum_{k=1}^{n} f(k)-\int_{1}^{n} f(x) \mathrm{d} x\right]\) 收敛．

\subsection{题目6}
六．（15 分）计算广义积分 \(\int_{1}^{+\infty} \frac{\{x\}}{x^{3}} \mathrm{~d} x\) ，其中 \(\{x\}\) 表示 \(x\) 的小数部分，例如 \(\{\pi\}=\pi-3\) ．

\subsection{题目7}
七．（15 分）设 \(\delta \in[0,+\infty)\) ，讨论函数

\[
f(x, y)= \begin{cases}\frac{\sin \left(x^{1+\delta} y\right)}{x^{2}+y^{2}}, & (x, y) \neq(0,0) \\ 0, & (x, y)=(0,0)\end{cases}
\]

的连续性，可微性．

\subsection{题目8}
八．（15 分）设 \(f(x)\) 在 \([0,+\infty)\) 上连续．

\[
f(t)=t+\frac{1}{4 \pi} \iiint_{x^{2}+y^{2}+z^{2} \leq t^{2}} \frac{f\left(\sqrt{x^{2}+y^{2}+z^{2}}\right)}{x^{2}+y^{2}+z^{2}} \mathrm{~d} x \mathrm{~d} y \mathrm{~d} z, t \in[0,+\infty) .
\]

求 \(f(t)\) ．

\subsection{题目9}
九．（15 分）求曲线积分

\[
I=\int_{L} \frac{(x-y) \mathrm{d} x+(x+4 y) \mathrm{d} y}{x^{2}+4 y^{2}}
\]

其中 \(L\) 从 \((1,0)\) 沿上半圆周 \(x^{2}+y^{2}=1\) 到（－1，0）．

\subsection{题目10}
十．（15 分）在空间直角坐标系中，已知方程 \(\cos (3 x+2 z)=h(4 y-z)\) ，且函数 \(h\) 在 \(\mathbb{R}\) 上可微，证明：上述方程表示空间中的一个柱面。
\end{document}
